\chapter{绪论}
\vspace{7pt}

\section{课题背景}

随着互联网服务、云计算平台和分布式系统的快速发展，服务器系统已经成为企业业务承载、数据存储、在线交易、智能运维和网络安全防护的重要基础设施。因此，如何从海量运行数据中及时发现系统异常状态，已经成为可靠性工程、智能运维和异常检测领域的重要研究问题\cite{chandola2009anomaly,he2016experience}。

服务器日志是系统运行过程中自动生成的重要文本数据，记录了服务器在不同时间点发生的事件、状态变化、错误信息、用户行为以及进程调用过程。因此，基于日志分析实现服务器异常状态检测，对于故障预警、风险控制和系统维护具有重要价值。\cite{oliner2007supercomputers,xu2009detecting}\cite{chandola2009anomaly,he2016experience}\cite{vaarandi2015logcluster,du2016spell,he2017drain,zhu2019tools}\cite{du2017deeplog,meng2019loganomaly,guo2021logbert}\cite{hochreiter1997lstm}\cite{vaswani2017attention}\cite{he2009imbalanced,saito2015precision}\cite{hasani2021liquid,chen2018neuralode}


\section{课题研究意义}

本课题围绕服务器日志异常检测问题，设计并实现一种基于日志分析与 Liquid 架构的服务器异常检测系统。系统从原始日志数据出发，依次完成日志预处理、日志解析、特征工程、模型训练、阈值优化和异常判别等步骤，最终实现对服务器异常状态的自动识别。

首先，本课题有助于提升服务器运维的自动化与智能化水平。通过构建自动化日志异常检测系统，可以在海量日志中快速发现潜在异常模式，减少人工排查成本，使运维人员能够更加及时地定位故障和处理风险。\cite{xu2009detecting,he2016experience}\cite{saito2015precision}


\section{国内外研究现状}

日志异常检测是智能运维领域的重要研究方向之一。随着系统规模不断扩大，国内外学者围绕日志数据集建设、日志解析、日志特征表示和异常检测模型等方面开展了大量研究。

在日志数据集与日志分析基础研究方面，Oliner 和 Stearley 对多类超级计算机系统日志进行了系统研究，指出系统日志能够反映大规模计算系统的运行状态和故障特征\cite{oliner2007supercomputers}。Xu 等通过挖掘控制台日志检测大规模系统问题，证明了从日志中提取事件模式和统计特征进行异常检测的可行性\cite{xu2009detecting}。\cite{zhu2023loghub}\cite{vaarandi2015logcluster}\cite{du2016spell}\cite{he2017drain}\cite{zhu2019tools}\cite{mikolov2013efficient,devlin2019bert}\cite{du2017deeplog}\cite{scholkopf2001estimating,liu2008isolation}\cite{meng2019loganomaly}\cite{guo2021logbert}\cite{he2009imbalanced,saito2015precision}\cite{vaswani2017attention,guo2021logbert}\cite{chen2018neuralode}\cite{hasani2021liquid,hasani2022closedform}\cite{lechner2020neural}


\section{本文主要研究内容}

本文以服务器日志数据为研究对象，围绕日志异常检测问题，设计并实现一个基于日志分析与 Liquid 架构的服务器异常检测系统。本文的主要研究内容包括以下几个方面。

\subsection{日志数据预处理}

原始服务器日志通常存在格式不统一、字段冗余、噪声信息较多和数据缺失等问题。本文首先对日志数据库中的原始数据进行清洗与规范化处理，包括日志数据读取、无效字段过滤、时间戳统一转换、重复日志去除以及特殊字符清理等。

\subsection{基于规则与正则表达式的日志解析}

针对预处理后的非结构化日志文本，本文设计日志解析模块。该模块通过规则和正则表达式识别日志中的动态参数，例如 IP 地址、端口号、时间字段、进程编号、内存地址和数字变量等，并将其替换为统一占位符，从而提取代表核心语义的固定日志模板。\cite{vaarandi2015logcluster,he2017drain,zhu2019tools}

\subsection{面向连续时序的特征工程构建}

为适配 Liquid 模型的连续时间建模特点，本文设计联合特征构建方法。通过这种方式，模型输入同时包含局部统计分布、事件顺序和时间间隔信息，有助于更全面地描述服务器运行状态。

\subsection{基于 Liquid 架构的模型训练与动态阈值寻优}

针对传统离散时序模型难以充分利用非规则日志时间间隔的问题，本文采用 Liquid 架构构建异常检测核心模型。模型以日志窗口特征序列为输入，通过连续时间动态机制学习服务器正常运行状态下的状态演化规律，并输出异常评分或异常概率。\cite{he2009imbalanced,saito2015precision}

\subsection{异常检测系统设计、实现与实验验证}

在算法研究的基础上，本文实现一套端到端服务器日志异常检测系统。系统整合日志读取、数据清洗、模板解析、特征工程、模型训练、动态阈值选择、异常推断和结果输出等模块，实现数据自动化流转和异常状态识别。

\section{本文主要贡献与创新点}

针对现有服务器日志异常检测方法在非规则时序建模、类别不平衡处理和工程落地方面的不足，本文开展了模型设计与系统实现工作，主要贡献与创新点如下。

\begin{enumerate}
    \item \textbf{探索了 Liquid 架构在服务器日志异常检测中的应用。} 本文将基于连续时间动态系统的 Liquid 架构引入日志异常检测任务，利用其对时间间隔和输入变化的动态响应能力，刻画服务器日志在非规则采样条件下的状态演化过程，为离散日志序列建模提供了新的实现思路。
    
    \item \textbf{提出了一种融合统计频次与事件序列的联合特征构建方法。} 针对单一特征表达能力有限的问题，本文在时间窗口机制下同时构建事件频次分布特征和事件顺序编码特征，使模型既能够感知局部事件密度变化，也能够保留日志事件之间的时序依赖关系。
    
    \item \textbf{设计了面向类别不平衡问题的动态阈值寻优策略。} 针对真实服务器日志中异常样本占比较低的问题，本文基于验证集 Precision-Recall 曲线选择异常判定阈值，避免简单采用固定阈值造成误报或漏报偏高，从而提升模型在不平衡数据场景下的综合检测效果。
    
    \item \textbf{实现并验证了一套端到端自动化日志异常检测系统。} 本文打通了从日志读取、清洗、解析、特征构建，到模型训练、异常推断和结果输出的完整流程，并在公开日志数据集上进行实验验证，证明了系统在服务器异常检测任务中的可行性和工程应用价值。
\end{enumerate}

\section{本文组织结构}

本文共分为六章，各章节内容安排如下。

\textbf{第一章 绪论。} 主要阐述服务器日志异常检测的课题背景与研究意义，梳理国内外在日志解析、特征表示、异常检测模型和连续时间建模方面的研究现状，指出现有方法的不足，并概述本文的主要研究内容、贡献与论文组织结构。

\textbf{第二章 相关理论与关键技术。} 介绍服务器日志的数据特性与分析难点，回顾日志预处理、日志解析、特征工程和异常检测的常用方法；重点阐述深度学习时序模型、Liquid 架构以及 Liquid Time-Constant Networks 的基本思想，为后续系统设计与模型实现奠定理论基础。

\textbf{第三章 系统需求分析与总体设计。} 从实际运维应用场景出发，对系统进行功能需求和非功能需求分析，论证系统建设可行性；随后提出系统总体架构，明确数据层、处理层、模型层和输出层的模块划分，并设计核心数据流转过程和数据库结构。

\textbf{第四章 系统实现。} 详细描述系统各功能模块的具体实现过程，包括日志读取与清洗、正则表达式驱动的日志解析、基于时间窗口的特征工程构建、Liquid 异常检测模型的训练与推断，以及动态阈值判定和异常检测结果输出逻辑。

\textbf{第五章 系统测试与实验分析。} 基于公开服务器日志数据集对本文系统进行功能测试和性能评估，展示模型训练过程、异常评分变化和检测结果，并使用准确率、精确率、召回率、F1 值等指标与相关基线方法进行对比分析；同时对输入序列长度、时间窗口大小和判定阈值等关键因素进行讨论。

\textbf{第六章 总结与展望。} 对全文研究工作进行总结，归纳本文取得的主要成果；同时分析当前系统在复杂日志泛化、深层语义利用和流式实时处理方面存在的不足，并提出后续改进方向和研究展望。